% Generated by `taskgen course examples` from dlp.elgamal (seed 5).
% Do not edit: change examples.yaml and regenerate.
\begin{example}[шифрование Эль-Гамаля]\label{ex:elgamal-toy}
Открытые параметры системы Эль-Гамаля: $p = 59$, $g = 50$; секретный ключ Алисы $a = 51$. Боб шифрует для Алисы сообщение $m = 2$ со случайным (одноразовым) показателем $k = 44$. \par (а) Найти открытый ключ Алисы $A$ и шифртекст $(c_1, c_2)$. \par (б) Расшифровать $(c_1, c_2)$ так, как это делает Алиса: обратный элемент найти расширенным алгоритмом Евклида и проверить по малой теореме Ферма.

\par\noindent\textit{Решение.} Степени вычисляем бинарным алгоритмом: биты показателя слева направо, $r \gets r^2$, при бите $1$ ещё $r \gets r \cdot (\text{основание})$; все вычисления по модулю $p = 59$.\par (а) $A = g^a$, $c_1 = g^k$, $c_2 = m \cdot A^k$:\[A = 50^{51}:\quad\begin{array}{c|rrrrrr}\text{бит} & 1 & 1 & 0 & 0 & 1 & 1 \\ \hline r^2 &  & 22 & 28 & 17 & 53 & 25 \\ r & 50 & 38 & 28 & 17 & 54 & 11\end{array}\qquad A \equiv 11.\]\[c_1 = 50^{44}:\quad\begin{array}{c|rrrrrr}\text{бит} & 1 & 0 & 1 & 1 & 0 & 0 \\ \hline r^2 &  & 22 & 12 & 41 & 48 & 3 \\ r & 50 & 22 & 10 & 44 & 48 & 3\end{array}\qquad c_1 \equiv 3.\]\[s = 11^{44}:\quad\begin{array}{c|rrrrrr}\text{бит} & 1 & 0 & 1 & 1 & 0 & 0 \\ \hline r^2 &  & 3 & 9 & 7 & 29 & 15 \\ r & 11 & 3 & 40 & 18 & 29 & 15\end{array}\qquad s \equiv 15.\]$c_2 = m s = 2 \cdot 15 = 30 \equiv 30$. Шифртекст $(c_1, c_2) = (3, 30)$.\par (б) Алиса знает $a$ и вычисляет маску $s = c_1^a = g^{ka} = A^k$:\[s = 3^{51}:\quad\begin{array}{c|rrrrrr}\text{бит} & 1 & 1 & 0 & 0 & 1 & 1 \\ \hline r^2 &  & 9 & 21 & 28 & 17 & 5 \\ r & 3 & 27 & 21 & 28 & 51 & 15\end{array}\qquad s \equiv 15.\] Обратный к $15$ по модулю $59$ (расширенный алгоритм Евклида для $(59, 15)$):
\[\begin{array}{r|r|r|r}
r_i & q_i & u_i & v_i \\ \hline
59 &  & 1 & 0 \\
15 &  & 0 & 1 \\
14 & 3 & 1 & -3 \\
1 & 1 & -1 & 4 \\
0 & 14 & 15 & -59 \\
\end{array}\qquad r_i = u_i \cdot 59 + v_i \cdot 15\]

$15^{-1} \equiv 4\ (\textrm{mod}\ 59)$.
Проверка по малой теореме Ферма: $s^{-1} = c_1^{-a} = c_1^{p - 1 - a} = 3^{7} \equiv 4$. Тогда $m = c_2 s^{-1} = 30 \cdot 4 = 120 \equiv 2$.\par $\ord(g) = 58$ ($g$ --- первообразный корень): маска $s = A^k$ пробегает всю группу $(\mathbb{Z}/59\mathbb{Z})^*$.\[A = 11,\quad (c_1, c_2) = (3, 30),\quad m = 2.\]
\end{example}
